\chapter{Spécification et analyse des besoins}

\section*{Introduction}
Ce chapitre transforme le cahier des charges en un périmètre vérifiable. Il identifie
les acteurs humains et externes, formalise les besoins fonctionnels et non
fonctionnels, puis présente le Product Backlog révisé et sa planification en Sprints.
Les interactions, scénarios détaillés et diagrammes associés seront étudiés dans le
chapitre de conception afin de séparer clairement la spécification de la modélisation
de la solution.

\section{Spécification des besoins}

\subsection{Identification des acteurs}

Un acteur représente une entité extérieure au système qui poursuit un objectif au
moyen de celui-ci. Une base de données, un worker de file ou le scheduler Laravel ne
sont donc pas des acteurs : ils appartiennent à l'implémentation interne. En revanche,
une borne, Google OAuth ou un prestataire de paiement sont des acteurs secondaires,
car ils échangent avec la plateforme par une interface clairement délimitée.

La notion d'\textbf{utilisateur authentifié} est utilisée comme acteur abstrait pour
regrouper les capacités communes : connexion, profil, paramètres personnels,
onboarding et notifications. Les cinq rôles authentifiés spécialisés restent le
Super Administrateur, l'Administrateur, l'Opérateur, le Technicien et le Client.

\begin{longtable}{|C{3.15cm}|C{2.15cm}|p{8.15cm}|}
  \caption{Acteurs humains de la plateforme}
  \label{tab:human-actors}\\
  \hline
  \textbf{Acteur} & \textbf{Catégorie} & \textbf{Responsabilités et périmètre}\\
  \hline
  \endfirsthead
  \hline
  \textbf{Acteur} & \textbf{Catégorie} & \textbf{Responsabilités et périmètre}\\
  \hline
  \endhead
  Visiteur & Principal &
  Consulte la présentation publique, les informations publiques des bornes, contacte
  la plateforme, crée un compte client ou dépose une demande de démonstration.\\\hline
  Client & Principal &
  Utilisateur global sans organisation fixe. Il peut utiliser les bornes de plusieurs
  organisations, gérer ses abonnements, lancer ses propres sessions, payer et
  consulter uniquement ses données.\\\hline
  Opérateur & Principal &
  Employé d'une seule organisation. Il met en service les stations, supervise OCPP,
  traite les alertes, exécute les commandes distantes, assigne les interventions et
  planifie les maintenances dans son périmètre.\\\hline
  Technicien & Principal &
  Employé d'une seule organisation. Il consulte les actifs en lecture, traite
  uniquement les interventions et maintenances qui lui sont assignées, puis produit
  les preuves et rapports techniques.\\\hline
  Administrateur & Principal &
  Responsable d'une seule organisation. Il gère ses opérateurs et techniciens, ses
  clients observés, ses actifs, tarifs, documents, rapports et son abonnement
  commercial, sans administrer la plateforme globale.\\\hline
  Super Administrateur & Principal &
  Gère les organisations et leurs administrateurs, les demandes de démonstration,
  les offres SaaS, l'audit, les rôles, les paramètres globaux et l'état des
  intégrations.\\
  \hline
\end{longtable}

\begin{longtable}{|C{3.15cm}|C{2.15cm}|p{8.15cm}|}
  \caption{Acteurs secondaires et systèmes externes}
  \label{tab:external-actors}\\
  \hline
  \textbf{Acteur} & \textbf{Nature} & \textbf{Interaction avec \ProjectName{}}\\
  \hline
  \endfirsthead
  \hline
  \textbf{Acteur} & \textbf{Nature} & \textbf{Interaction avec \ProjectName{}}\\
  \hline
  \endhead
  Borne ou simulateur OCPP & Système externe &
  Ouvre une connexion WebSocket authentifiée, émet les messages OCPP, reçoit les
  commandes distantes et retourne leur résultat.\\\hline
  Google OAuth & Fournisseur d'identité &
  Authentifie un compte Google vérifié sans transmettre son mot de passe à la
  plateforme.\\\hline
  Service de paiement & Service externe &
  Autorise, capture ou libère un montant et envoie des webhooks signés. Dans le MVP,
  ce service est simulé mais reste hors du domaine métier.\\\hline
  Service email & Service externe &
  Distribue les vérifications, invitations, réinitialisations et notifications
  transactionnelles placées dans la file.\\\hline
  Service cartographique & Service externe &
  Fournit les tuiles et données nécessaires à l'affichage de la carte et à
  l'ouverture d'un itinéraire.\\
  \hline
\end{longtable}

\begin{scopebox}[Clarification du rôle financier]
  La \emph{finance} n'est pas modélisée comme un acteur autonome. Les décisions
  commerciales sont prises par le Super Administrateur ou l'Administrateur selon
  leur périmètre. Le seul acteur externe du paiement est le prestataire technique qui
  traite une transaction. Cette distinction évite de confondre une fonction métier,
  un rôle humain et un système tiers.
\end{scopebox}

\subsection{Besoins fonctionnels}

Les besoins fonctionnels décrivent ce que le système doit permettre de réaliser. Ils
sont regroupés par acteur, mais les contrôles d'autorisation restent appliqués côté
API même lorsqu'une action n'est pas affichée dans l'interface.

\subsubsection{Fonctions communes aux utilisateurs authentifiés}

\begin{longtable}{|C{1.35cm}|p{12.3cm}|}
  \hline
  \textbf{Code} & \textbf{Besoin fonctionnel}\\
  \hline
  \endfirsthead
  \hline
  \textbf{Code} & \textbf{Besoin fonctionnel}\\
  \hline
  \endhead
  BF-01 & Se connecter par email et mot de passe ou, lorsque le compte le permet,
  par Google OAuth ; se déconnecter et récupérer son accès.\\\hline
  BF-02 & Consulter et modifier les données personnelles autorisées, l'adresse,
  l'avatar et les liens professionnels pertinents.\\\hline
  BF-03 & Configurer le fuseau horaire, le rayon de recherche, les préférences de
  notification et le mot de passe lorsqu'un mot de passe local existe.\\\hline
  BF-04 & Suivre un onboarding court adapté au rôle, puis retrouver des indications
  contextuelles sans être bloqué dans une visite générique.\\\hline
  BF-05 & Consulter les notifications personnelles, leur contexte, leur statut lu ou
  non lu et marquer un ensemble cohérent comme lu.\\\hline
  BF-06 & Utiliser une recherche globale limitée aux objets réellement accessibles.\\
  \hline
\end{longtable}

\subsubsection{Super Administrateur}

\begin{longtable}{|C{1.35cm}|p{12.3cm}|}
  \hline
  \textbf{Code} & \textbf{Besoin fonctionnel}\\
  \hline
  BF-07 & Qualifier, rejeter ou convertir une demande de démonstration en organisation
  d'essai et invitation d'administrateur.\\\hline
  BF-08 & Créer, consulter, modifier, activer ou suspendre les organisations et gérer
  leurs administrateurs.\\\hline
  BF-09 & Consulter les indicateurs globaux, l'audit filtrable et la matrice des rôles
  et permissions.\\\hline
  BF-10 & Consulter l'état des intégrations sans exposer leurs secrets et modifier
  uniquement les paramètres globaux autorisés.\\\hline
  BF-11 & Gérer les offres SaaS, essais, capacités, abonnements et factures des
  organisations, avec historique des transitions.\\\hline
  BF-12 & Produire les analyses et exports globaux de la plateforme.\\
  \hline
\end{longtable}

\subsubsection{Administrateur d'organisation}

\begin{longtable}{|C{1.35cm}|p{12.3cm}|}
  \hline
  \textbf{Code} & \textbf{Besoin fonctionnel}\\
  \hline
  BF-13 & Inviter, modifier, désactiver ou supprimer les opérateurs et techniciens de
  sa seule organisation, sans connaître leur mot de passe.\\\hline
  BF-14 & Consulter les clients ayant réellement utilisé les bornes de l'organisation,
  sans modifier leur compte global ni leurs activités chez d'autres opérateurs.\\\hline
  BF-15 & Administrer les stations, connecteurs, commandes, alertes, interventions,
  maintenances et documents de son organisation.\\\hline
  BF-16 & Définir les tarifs, plans clients et affectations applicables aux stations
  ou connecteurs de son organisation.\\\hline
  BF-17 & Consulter les indicateurs métier, produire et échanger des rapports et
  exporter les données autorisées.\\\hline
  BF-18 & Suivre l'offre commerciale, les quotas, les échéances et factures de son
  organisation et demander un changement d'offre.\\
  \hline
\end{longtable}

\subsubsection{Opérateur}

\begin{longtable}{|C{1.35cm}|p{12.3cm}|}
  \hline
  \textbf{Code} & \textbf{Besoin fonctionnel}\\
  \hline
  BF-19 & Mettre en service une station et ses connecteurs de manière atomique, la
  positionner sur la carte et préparer sa liaison à une borne physique ou simulée.\\\hline
  BF-20 & Consulter les stations en vue carte, liste ou tableau avec disponibilité
  calculée, connexion, dernière preuve de vie et alertes.\\\hline
  BF-21 & Examiner une station, sa télémétrie vérifiée, ses connecteurs, événements,
  commandes, sessions, maintenances, alertes et documents.\\\hline
  BF-22 & Exécuter les commandes OCPP autorisées : redémarrage logiciel,
  déverrouillage et changement de disponibilité pour maintenance.\\\hline
  BF-23 & Filtrer, acquitter, assigner, suivre ou annuler une alerte dans le respect
  de son cycle de vie et de son SLA.\\\hline
  BF-24 & Créer et suivre les interventions et maintenances préventives ou correctives,
  y compris leur récurrence et leur calendrier.\\\hline
  BF-25 & Superviser les sessions et paiements de l'organisation sans agir à la place
  d'un client sur ses données privées.\\\hline
  BF-26 & Consulter et échanger des rapports opérationnels adaptés à son rôle.\\
  \hline
\end{longtable}

\subsubsection{Technicien}

\begin{longtable}{|C{1.35cm}|p{12.3cm}|}
  \hline
  \textbf{Code} & \textbf{Besoin fonctionnel}\\
  \hline
  BF-27 & Consulter en lecture les stations, connecteurs, cartes et informations
  techniques de son organisation.\\\hline
  BF-28 & Consulter les alertes, interventions et maintenances qui lui sont assignées,
  avec priorité, échéance et historique.\\\hline
  BF-29 & Faire progresser une intervention selon les transitions autorisées, ajouter
  des notes et téléverser des preuves avant/après privées.\\\hline
  BF-30 & Produire un rapport final comprenant diagnostic, actions, pièces, contrôles
  terrain et résultat, puis demander un suivi lorsque nécessaire.\\\hline
  BF-31 & Consulter ses propres indicateurs de terrain et échanger des rapports avec
  les destinataires autorisés de son organisation.\\
  \hline
\end{longtable}

\subsubsection{Client}

\begin{longtable}{|C{1.35cm}|p{12.3cm}|}
  \hline
  \textbf{Code} & \textbf{Besoin fonctionnel}\\
  \hline
  BF-32 & Rechercher les stations publiques des organisations actives, utiliser la
  géolocalisation selon un rayon choisi et filtrer les résultats.\\\hline
  BF-33 & Consulter les connecteurs et le tarif effectif, copier les coordonnées ou
  ouvrir un itinéraire vers la station sélectionnée.\\\hline
  BF-34 & Démarrer une recharge guidée : choisir station et connecteur, brancher le
  câble, attendre sa détection OCPP, autoriser le paiement puis lancer la commande.\\\hline
  BF-35 & Suivre en temps réel la durée, l'énergie, la puissance et le coût provisoire,
  puis arrêter la session manuellement ou par une limite configurée.\\\hline
  BF-36 & Consulter l'historique multi-organisations de ses sessions et paiements,
  prévisualiser ou télécharger un reçu.\\\hline
  BF-37 & Consulter les plans proposés par les organisations et gérer plusieurs
  abonnements indépendants avec leurs factures.\\\hline
  BF-38 & Utiliser un identifiant OCPP virtuel ; préparer l'évolution vers une carte
  RFID ou la lecture d'un QR code physique sans exposer de secret.\\
  \hline
\end{longtable}

\subsubsection{Visiteur et traitements automatiques}

\begin{longtable}{|C{1.35cm}|p{12.3cm}|}
  \hline
  \textbf{Code} & \textbf{Besoin fonctionnel}\\
  \hline
  BF-39 & Présenter la solution et les moyens de contact, permettre l'inscription
  client et la demande de démonstration sans exposer les données protégées.\\\hline
  BF-40 & Recevoir, authentifier et rendre idempotents les événements OCPP, puis
  projeter les états métier et leurs transitions.\\\hline
  BF-41 & Détecter les absences de communication, défauts, échéances SLA ou
  maintenances proches et produire les alertes et notifications correspondantes.\\\hline
  BF-42 & Traiter en tâche de fond les emails, occurrences récurrentes, contrôles de
  SLA, exports volumineux et événements de paiement.\\\hline
  BF-43 & Vérifier les signatures et l'idempotence des webhooks de paiement, puis
  réconcilier le paiement, la session et le reçu.\\
  \hline
\end{longtable}

\subsection{Besoins non fonctionnels}

Les objectifs du cahier des charges sont transformés en exigences vérifiables. Les
valeurs de production qui ne peuvent pas être garanties par un environnement local
sont explicitement présentées comme des objectifs à valider lors du déploiement.

\begin{longtable}{|C{1.45cm}|C{3.0cm}|p{9.2cm}|}
  \caption{Exigences non fonctionnelles}
  \label{tab:non-functional-requirements}\\
  \hline
  \textbf{Code} & \textbf{Qualité} & \textbf{Exigence et critère de vérification}\\
  \hline
  \endfirsthead
  \hline
  \textbf{Code} & \textbf{Qualité} & \textbf{Exigence et critère de vérification}\\
  \hline
  \endhead
  ENF-01 & Performance &
  Sous charge nominale et avec pagination, une page courante doit devenir utilisable
  en moins de deux secondes ; les opérations longues sont asynchrones et signalent
  leur progression.\\\hline
  ENF-02 & Temps réel &
  Un événement OCPP accepté doit mettre à jour sa projection et être diffusé sans
  rechargement manuel. Une perte de communication est détectée après 90 secondes
  sans message, conformément à la règle métier retenue.\\\hline
  ENF-03 & Disponibilité du service &
  La cible de production est de 99,9\,\% hors maintenance planifiée. Cette cible doit
  être mesurée par supervision après déploiement ; elle n'est pas déclarée atteinte
  par le seul MVP local.\\\hline
  ENF-04 & Sécurité des accès &
  Toute route protégée utilise une session Sanctum, la protection CSRF, RBAC et une
  policy sur la ressource. Un identifiant deviné d'une autre organisation ne doit
  retourner aucune donnée.\\\hline
  ENF-05 & Isolation multi-tenant &
  Un administrateur, opérateur ou technicien appartient à exactement une organisation.
  Toutes les lectures, écritures, diffusions et exports sont limités à ce périmètre.\\\hline
  ENF-06 & Protection des secrets &
  Aucun secret OAuth, email, paiement ou OCPP n'est inclus dans le frontend, les logs
  applicatifs ou une réponse API. Les mots de passe et jetons persistés sont condensés
  ou chiffrés selon leur usage.\\\hline
  ENF-07 & Intégrité et idempotence &
  Une station et ses connecteurs sont créés atomiquement. Les événements OCPP,
  webhooks, commandes et tâches récurrentes supportent la répétition sans doublon
  métier.\\\hline
  ENF-08 & Traçabilité &
  Les changements sensibles conservent acteur, date, ancienne et nouvelle valeur ou
  événement métier. Les rapports finaux et reçus restent reproductibles.\\\hline
  ENF-09 & Maintenabilité &
  Les contrôleurs délèguent les règles aux services et policies ; les contrats
  frontend sont centralisés ; les modules critiques sont couverts par des tests
  automatisés et documentés.\\\hline
  ENF-10 & Extensibilité &
  Les fournisseurs de paiement, email, cartographie et identité sont intégrés derrière
  des contrats ou configurations remplaçables. OCPP reste séparé de l'API métier.\\\hline
  ENF-11 & Ergonomie &
  Les parcours critiques utilisent des étapes courtes, des états de chargement et des
  erreurs actionnables. L'interface reste cohérente entre rôles sans rendre leurs
  tableaux de bord identiques.\\\hline
  ENF-12 & Responsive et compatibilité &
  Les parcours principaux sont utilisables sur ordinateur et tablette, et les parcours
  client essentiels sur mobile. Les dernières versions de Chrome, Edge et Firefox
  sont ciblées.\\\hline
  ENF-13 & Confidentialité &
  La position courante du client reste dans l'état transitoire du navigateur et n'est
  pas persistée. Les pièces d'intervention privées sont servies après autorisation.\\\hline
  ENF-14 & Portabilité &
  L'environnement peut être reproduit avec Git, PHP, Node.js, PostgreSQL et Docker
  sous WSL2 pour les services simulés, sans dépendre d'un poste unique.\\\hline
  ENF-15 & Interopérabilité &
  Les échanges métier utilisent une API REST JSON documentable par OpenAPI ; les
  stations utilisent OCPP 1.6 JSON sur WebSocket et le paiement suit un contrat
  fournisseur indépendant.\\
  \hline
\end{longtable}

\section{Planification du projet}

\subsection{Méthode de priorisation et d'estimation}

Deux informations distinctes sont associées à chaque user story :

\begin{itemize}[leftmargin=7mm]
  \item la \textbf{priorité} exprime la valeur métier et l'urgence. Les niveaux Haute,
  Moyenne et Faible correspondent respectivement à une adaptation des catégories
  \emph{Must}, \emph{Should} et \emph{Could} de MoSCoW ;
  \item les \textbf{story points} estiment l'effort relatif, l'incertitude et le risque
  selon la suite de Fibonacci : 1, 2, 3, 5, 8 et 13.
\end{itemize}

Les points ne représentent donc ni des jours ni une priorité. Une fonctionnalité
hautement prioritaire peut recevoir peu de points si elle est simple, tandis qu'une
évolution facultative comme le firmware peut recevoir 13 points en raison de son
risque technique.

\subsection{Product Backlog révisé}

Le backlog initial de 30 user stories a été enrichi à mesure que les relations
multi-organisations, la disponibilité calculée et les parcours réels ont été
clarifiés. La gestion des véhicules, retirée du produit après validation UX, est
remplacée par un onboarding adapté au rôle. Deux user stories distinguent désormais
la gestion commerciale globale de la consultation de l'abonnement par
l'administrateur. Le backlog révisé contient 42 user stories pour 301 points.

\begingroup
\footnotesize
\setlength{\tabcolsep}{3.5pt}
\begin{longtable}{|C{1.25cm}|p{10.25cm}|C{1.65cm}|C{1.0cm}|}
  \caption{Product Backlog de \ProjectName{}}
  \label{tab:product-backlog}\\
  \hline
  \textbf{ID} & \textbf{User Story} & \textbf{Priorité} & \textbf{Pts}\\
  \hline
  \endfirsthead
  \hline
  \textbf{ID} & \textbf{User Story} & \textbf{Priorité} & \textbf{Pts}\\
  \hline
  \endhead
  US-01 & En tant que visiteur, je veux créer un compte client ou me connecter de
  manière sécurisée par email ou Google. & Haute & 8\\\hline
  US-02 & En tant qu'utilisateur authentifié, je veux consulter et mettre à jour mon
  profil ainsi que les paramètres de sécurité de mon compte. & Moyenne & 3\\\hline
  US-03 & En tant que super administrateur, je veux gérer les organisations, leurs
  administrateurs et les accès globaux de la plateforme. & Haute & 8\\\hline
  US-04 & En tant qu'administrateur, je veux créer, modifier et désactiver les comptes
  des opérateurs et techniciens de mon organisation. & Haute & 5\\\hline
  US-05 & En tant que super administrateur, je veux consulter l'audit global, l'état
  des intégrations et les statistiques d'utilisation de la plateforme. & Moyenne & 5\\\hline
  US-06 & En tant qu'administrateur, je veux consulter les clients ayant utilisé les
  bornes de mon organisation avec leurs sessions et statistiques. & Moyenne & 5\\\hline
  US-07 & En tant qu'opérateur, je veux ajouter, modifier et désactiver les bornes de
  mon organisation avec leurs informations techniques et géographiques. & Haute & 8\\\hline
  US-08 & En tant qu'opérateur, je veux gérer les connecteurs associés aux bornes de
  mon organisation avec leur type, puissance et état. & Haute & 5\\\hline
  US-09 & En tant qu'opérateur, je veux consulter la liste des bornes avec leur
  disponibilité calculée, leur connexion et leur dernière activité. & Haute & 8\\\hline
  US-10 & En tant qu'opérateur, je veux visualiser les bornes sur une carte et ajouter
  une borne en sélectionnant son emplacement. & Haute & 8\\\hline
  US-11 & En tant que technicien, je veux consulter sans modification les bornes et
  leur emplacement sur la carte de mon organisation. & Haute & 5\\\hline
  US-12 & En tant que client, je veux rechercher et filtrer les bornes publiques selon
  leur disponibilité, connecteur et tarif. & Haute & 8\\\hline
  US-13 & En tant qu'opérateur, je veux consulter la fiche détaillée d'une borne avec
  ses connecteurs, sessions, historique et alertes. & Haute & 8\\\hline
  US-14 & En tant que plateforme, je veux recevoir et valider automatiquement les
  messages OCPP 1.6 JSON des bornes physiques ou simulées. & Haute & 13\\\hline
  US-15 & En tant que plateforme, je veux calculer l'état des bornes à partir des
  événements OCPP, transactions, heartbeats et délais de non-réponse. & Haute & 13\\\hline
  US-16 & En tant qu'opérateur, je veux superviser en temps réel les états disponible,
  occupée, non disponible, maintenance, déconnectée ou en défaut. & Haute & 13\\\hline
  US-17 & En tant qu'opérateur, je veux consulter un tableau de bord avec disponibilité,
  énergie, sessions actives et taux d'indisponibilité. & Haute & 8\\\hline
  US-18 & En tant qu'administrateur, je veux consulter le tableau de bord de mon
  organisation avec utilisateurs, clients, revenus, régions et bornes. & Haute & 8\\\hline
  US-19 & En tant qu'opérateur, je veux recevoir des alertes lorsqu'une borne est
  déconnectée, en panne, en surchauffe ou en erreur de communication. & Haute & 8\\\hline
  US-20 & En tant qu'opérateur, je veux filtrer, suivre et assigner les alertes selon
  leur gravité, statut et borne. & Moyenne & 5\\\hline
  US-21 & En tant que technicien, je veux consulter les alertes et interventions qui
  me sont assignées. & Haute & 5\\\hline
  US-22 & En tant que technicien, je veux créer un rapport d'intervention avec
  diagnostic, actions, pièces et photos. & Moyenne & 8\\\hline
  US-23 & En tant qu'opérateur, je veux planifier et suivre les maintenances
  préventives et correctives. & Moyenne & 8\\\hline
  US-24 & En tant que client, je veux démarrer une session sur une borne disponible,
  quelle que soit son organisation. & Haute & 8\\\hline
  US-25 & En tant que client, je veux arrêter ma session et consulter le résumé de
  durée, énergie et coût. & Haute & 5\\\hline
  US-26 & En tant qu'utilisateur autorisé, je veux consulter l'historique des sessions
  de mon périmètre avec durée, énergie, coût et borne. & Moyenne & 5\\\hline
  US-27 & En tant que nouvel utilisateur, je veux un onboarding adapté à mon rôle afin
  de configurer mon espace et atteindre une première action utile. & Moyenne & 5\\\hline
  US-28 & En tant que client, je veux utiliser une carte RFID ou un QR code pour
  m'identifier et lancer une recharge. & Moyenne & 8\\\hline
  US-29 & En tant que client, je veux payer mes sessions, suivre le paiement et
  prévisualiser ou télécharger mes reçus. & Moyenne & 8\\\hline
  US-30 & En tant que plateforme, je veux traiter les paiements au moyen d'un adaptateur
  simulé et remplaçable. & Moyenne & 8\\\hline
  US-31 & En tant qu'administrateur, je veux configurer les tarifs selon le kWh, la
  durée ou les plages horaires et les affecter aux actifs. & Moyenne & 8\\\hline
  US-32 & En tant que client, je veux choisir les plans des organisations et gérer
  simultanément plusieurs abonnements. & Moyenne & 8\\\hline
  US-33 & En tant qu'utilisateur autorisé, je veux consulter des analyses adaptées à
  mon rôle et échanger des rapports opérationnels. & Moyenne & 8\\\hline
  US-34 & En tant qu'utilisateur autorisé, je veux exporter les données et rapports de
  mon périmètre en PDF, JSON ou CSV. & Faible & 5\\\hline
  US-35 & En tant qu'administrateur, je veux gérer les paramètres locaux de mon
  organisation : TVA, logo, langue, fuseau et notifications. & Faible & 5\\\hline
  US-36 & En tant que super administrateur, je veux configurer les intégrations
  globales : paiement, notifications, cartographie, OAuth2, OCPP et sécurité. & Faible & 8\\\hline
  US-37 & En tant qu'administrateur, je veux gérer les contrats, garanties, notices et
  plans techniques liés aux bornes. & Faible & 5\\\hline
  US-38 & En tant qu'administrateur, je veux gérer les mises à jour firmware des bornes
  à distance. & Faible & 13\\\hline
  US-39 & En tant que visiteur, je veux consulter la présentation, contacter la
  plateforme et envoyer une demande de démonstration. & Moyenne & 3\\\hline
  US-40 & En tant qu'utilisateur authentifié, je veux recevoir des notifications
  in-app et email et gérer mes préférences. & Moyenne & 5\\\hline
  US-41 & En tant que super administrateur, je veux gérer les offres SaaS, essais,
  abonnements et factures des organisations. & Haute & 8\\\hline
  US-42 & En tant qu'administrateur, je veux suivre l'abonnement commercial, les
  limites et factures de mon organisation et demander un changement d'offre. & Moyenne & 5\\
  \hline
  \multicolumn{2}{r}{\textbf{Total}} &
  \multicolumn{2}{c}{\textbf{301 points}}\\
  \hline
\end{longtable}
\endgroup

\begin{table}[H]
  \centering
  \caption{Répartition du backlog par priorité}
  \label{tab:backlog-priority-summary}
  \begin{tabular}{@{}lrrr@{}}
    \toprule
    \textbf{Priorité} & \textbf{User stories} & \textbf{Points} &
    \textbf{Part des points}\\
    \midrule
    Haute & 20 & 160 & 53,2\,\%\\
    Moyenne & 17 & 105 & 34,9\,\%\\
    Faible & 5 & 36 & 12,0\,\%\\
    \midrule
    \textbf{Total} & \textbf{42} & \textbf{301} & \textbf{100\,\%}\\
    \bottomrule
  \end{tabular}
\end{table}

\subsection{Planification des Sprints}

Le Product Backlog est réalisé au moyen de Sprints courts, chacun associé à un
objectif vérifiable. La cadence prévue est hebdomadaire ; le dernier Sprint est
légèrement raccourci afin de respecter la date de fin du stage. Les périodes du
tableau~\ref{tab:sprint-backlogs} représentent la planification de référence. Le
contenu détaillé d'un Sprint peut être ajusté lors du Sprint Planning ou réordonné
après une Sprint Review sans modifier l'objectif global du produit.

Les story points ne sont pas convertis en heures. Ils servent à comparer l'effort et
le risque des user stories. La capacité est appréciée progressivement à partir des
Sprints précédents, ce qui évite de présenter une précision artificielle au début du
stage.

\begingroup
\small
\setlength{\tabcolsep}{2.5pt}
\begin{longtable}{|C{1.4cm}|C{2.25cm}|C{2.6cm}|p{3.3cm}|p{4.4cm}|}
  \caption{Synthèse des Sprint Backlogs}
  \label{tab:sprint-backlogs}\\
  \hline
  \textbf{Sprint} & \textbf{Période prévue} & \textbf{Objectif} &
  \textbf{User stories principales} & \textbf{Résultat inspecté en revue}\\
  \hline
  \endfirsthead
  \hline
  \textbf{Sprint} & \textbf{Période prévue} & \textbf{Objectif} &
  \textbf{User stories principales} & \textbf{Résultat inspecté en revue}\\
  \hline
  \endhead
  S01 & 1--7 juillet &
  Cadrer le produit et valider l'expérience cible &
  US-27, US-39 et travaux de conception &
  Backlog initial, acteurs, cas d'utilisation global, architecture de référence et
  prototype navigable.\\\hline
  S02 & 8--14 juillet &
  Sécuriser l'entrée et l'organisation des comptes &
  US-01 à US-06 &
  Authentification, création de compte, invitation, activation, rôles et isolation
  d'une organisation.\\\hline
  S03 & 15--21 juillet &
  Constituer et localiser le réseau de recharge &
  US-07 à US-13, US-37 &
  Stations, connecteurs, cartographie, fiche détaillée et documents accessibles selon
  le rôle.\\\hline
  S04 & 22--28 juillet &
  Produire une disponibilité issue du terrain &
  US-14 à US-16, US-19 &
  Connexion d'une station simulée, réception des événements, projection des états,
  historique et alerte de connectivité.\\\hline
  S05 & 29 juillet--4 août &
  Réaliser une recharge complète &
  US-24 à US-30 &
  Sélection d'une station, connexion, autorisation, mesures, arrêt, paiement simulé
  et reçu.\\\hline
  S06 & 5--11 août &
  Traiter un incident jusqu'à sa résolution &
  US-20 à US-23, US-40 &
  Alerte qualifiée, intervention assignée, maintenance, preuves de travail et
  notifications.\\\hline
  S07 & 12--18 août &
  Structurer les offres et la tarification &
  US-31, US-32, US-41, US-42 &
  Tarifs de recharge, adhésions des clients, essais et abonnements commerciaux des
  organisations.\\\hline
  S08 & 19--25 août &
  Fournir les outils de pilotage &
  US-05, US-17, US-18, US-33 à US-36 &
  Tableaux de bord spécialisés, rapports, exports, audit, paramètres et état des
  intégrations.\\\hline
  S09 & 26--31 août &
  Stabiliser et préparer la recette &
  Tests transversaux et documentation &
  Corrections finales, sécurité, tests de recette, API documentée, manuels, rapport
  et préparation du déploiement.\\
  \hline
\end{longtable}
\endgroup

La gestion distante du firmware (US-38) reste volontairement dans le Product Backlog
sans être engagée dans un Sprint du MVP. Elle nécessite une validation matérielle et
sécuritaire qui dépasse le parcours principal du stage. Cette transparence évite de
présenter comme réalisée une évolution seulement envisagée.

\subsection{Matrice de traçabilité fonctionnelle}

La matrice suivante relie les grands groupes de besoins à leurs principaux contrats
API et écrans. Les routes détaillées et la spécification OpenAPI seront présentées
dans les livrables techniques.

\begingroup
\small
\setlength{\tabcolsep}{3pt}
\begin{longtable}{|p{2.9cm}|p{2.9cm}|p{4.2cm}|p{3.9cm}|}
  \caption{Traçabilité entre backlog, API et interfaces}
  \label{tab:functional-traceability}\\
  \hline
  \textbf{Domaine} & \textbf{User stories} & \textbf{Contrats API principaux} &
  \textbf{Interfaces principales}\\
  \hline
  \endfirsthead
  \hline
  \textbf{Domaine} & \textbf{User stories} & \textbf{Contrats API principaux} &
  \textbf{Interfaces principales}\\
  \hline
  \endhead
  Identité et entrée & US-01--06, 27, 39 &
  Auth, profil, onboarding, utilisateurs, organisations, demandes de démo &
  Landing, connexion, activation, bienvenue, profil, utilisateurs, organisations\\\hline
  Stations et carte & US-07--13, 37 &
  Stations, commissioning, connecteurs, carte, télémétrie, documents &
  Stations, recherche de station, détail station, cartes par rôle\\\hline
  OCPP et disponibilité & US-14--16, 19 &
  API interne OCPP, commandes, transitions, alertes &
  Détail station, supervision, carte, centre d'alertes\\\hline
  Exploitation & US-20--23, 40 &
  Alertes, interventions, maintenances, notifications &
  Alertes, interventions, maintenance, calendrier, cloche de notifications\\\hline
  Recharge et paiement & US-24--30 &
  Tentatives, sessions, paiements, reçu, webhook fournisseur &
  Assistant de recharge, sessions, vue live, paiements et reçus\\\hline
  Tarification et commerce & US-31, 32, 41, 42 &
  Tarifs, plans, abonnements clients, portefeuille commercial, factures &
  Tarifs, abonnements, portefeuille commercial, facturation organisationnelle\\\hline
  Pilotage et rapports & US-05, 17--18,\newline 33--36 &
  Dashboards, reporting, rapports internes, exports, audit, intégrations, paramètres &
  Tableaux de bord par rôle, analyses, rapports, audit, intégrations, paramètres\\
  \hline
\end{longtable}
\endgroup


\section*{Conclusion}
La spécification distingue désormais clairement les responsabilités humaines, les
systèmes externes et les composants internes. Le backlog de 42 user stories couvre
l'entrée dans la plateforme, la supervision OCPP, l'exploitation terrain, la recharge,
le commerce et le pilotage, tout en conservant le firmware comme évolution à haut
risque. Les exigences non fonctionnelles rendent les objectifs de sécurité,
performance, isolation et traçabilité contrôlables. Le chapitre suivant traduira ces
besoins en architecture, cas d'utilisation détaillés, séquences et modèles de données.
